\magicSubsection{Variable Names}{sec:VariableNames}

Variable names $\sporkv$ are \emph{unique}:
each binder (procedure parameter, data or barrier allocation, window statement \texttt{$\sporkv$ = $\sporkw$ @ $\sporkspecialWindow$}, or loop variable)
introduces a name that is bound nowhere else in the program, so no name shadows another.

A binder binds its name for the remainder of the enclosing block (allocations, window statements) or for the loop body (loop variables);
assignment and reduction are not binders.
A variable is \emph{bound in} a statement $\sporks$ if its binder is inside $\sporks$, and \emph{free in} $\sporks$ if it is used in $\sporks$ but bound outside it
(e.g.\ procedure parameters, or enclosing loop variables and allocations).

Uniqueness is what lets the semantics stay simple:
substitution (Figure~\ref{fig:WindowSubSemantics}) never captures a variable,
and $\sporksigma$, which is keyed by name, needs no scoping:
loop variables and allocations simply stay in $\sporksigma$ after their scope ends, where nothing can observe them.

Rules that duplicate syntax must preserve uniqueness.
$\sporkalphaRename(\sporks)$ renames every variable bound in $\sporks$ to a fresh name, consistently throughout its scope,
and leaves variables free in $\sporks$ unchanged.
\textsc{C-ForIter} (Figure~\ref{fig:ControlSemantics}) renames each copy of the loop body,
so allocations in different iterations get distinct entries in $\sporksigma$,
while the loop variable itself (free in the body) is still read from $\sporksigma$.
Likewise, \sporkDInstrSub\ (Figure~\ref{fig:DataSemantics}) renames each inlined copy of $\sporkfBehavior$'s body
before substituting the instruction's parameters (free in the body) with the argument windows.
